\section{Conclusion}
Three different algorithms for the robot pick-and-place task have been tested. The classical algorithm is the most reliable. It also provides high control over the behavior. The issues with classification can be improved by retraining the NN with the suggestions that were mentioned in the previous section, and the issues with gripping can be improved by redesigning a gripper. The RL took a significant amount of time, and yet it is not able to complete the task. The control over the behavior is limited. The reward signal and hyperparameters need to be readjusted, and the model needs to be retrained. This is a lengthy and expensive process. The trajectory collection within VR with \nametwo was successful. The training took significantly less time, and the behavior is an improvement over the RL. The biggest issue with GAIL-IRL was the discriminator becoming strong early on, leading to premature convergence of the policy. These results show that manually programming the behavior can take less time and provide higher control over the behavior compared to the tested AI methods. Classical methods should be preferred in tasks where they can be implemented easily. 